\section{Edge-first coherent surfaces and a reusable span dual}
\label{sec:cocycle-lex-native}

Report 78 changes the optimization order. Earlier selectors minimize
normal-piece span $D$, then edge penalty $F$ within that minimum-span face.
The new selector minimizes $(F,D)$ lexicographically over all integer
vertex potentials in one coherent cocycle class. This reverses two genuine
objectives; minimizing $D$ first can exclude the best connected Euler
candidate. The native stage is optional because its current measurements
show additional cost without additional positive recognition coverage.

\subsection{Exact span bands and their real parameter}

Write $h_{t,i}$ for the four integral corner heights in tetrahedron $t$,
and $p_{v(t,i)}$ for global vertex potentials. Then
\begin{align*}
 D(p)&=\sum_t\bigl(\max_i(h_{t,i}+p_{v(t,i)})
                   -\min_i(h_{t,i}+p_{v(t,i)})\bigr),\\
 F(p)&=\sum_e w_e|c_e+p_{b(e)}-p_{a(e)}|,
 \qquad \chi=D-F/2.
\end{align*}
Here $w_e=\deg_{\mathrm{face}}(e)-2\ge0$ is authenticated from the actual
source. The coherent dual is cooriented and hence orientable. Let $D_*$
be the global minimum span with its independent primal/dual certificate.

The minimum-span anchor pairs make the sum of their height differences
equal to $D_*$: potential contributions cancel at each global vertex.
Consequently $D-D_*$ is the sum of $2T$ nonnegative integral defects from
the chosen lower and upper anchors. For radius $k$, pick $j$ positive
defects, one of three nonanchor attaining corners for each, and their
positive sizes with total at most $k$. This gives at most
\[
 H(T,k)=\sum_{j=0}^{\min(2T,k)}\binom{2T}{j}3^j\binom{k}{j}
           \le(1+6T)^k
\]
constant-span strata. Each has $O(T)$ difference constraints and permits
exact min-$F$ network optimization in polynomial bit time. Attainer ties
can repeat a potential in several descriptions without losing completeness.
When $k=O(\log n)$ and source size and coefficient bits are polynomial,
this supplied band can be searched in quasi-polynomial time.
There is no theorem that every unknot has a useful coherent disc in such
a band. An unrecorded streaming search also does not supply a portable
negative certificate for every stratum.

\subsection{Deleting zero-class coherent components}

\begin{proposition}[Positive core and excess budgets]
In an authenticated knot exterior, let a coherent surface $S$ represent
the positive primitive relative class and have $D=D_*+\delta$ pieces.
It contains a positive-class component $C$ with $D(C)\ge D_*$, and
\[
 D(S\setminus C)\le\delta,\qquad
 \#\pi_0(S)\le\delta+1.
\]
If there are $q$ negative-class components, then $2qD_*\le\delta$ and
the total zero-class mass is at most $\delta-2qD_*$. Also
$\chi(C)\ge\chi(S)-\delta$.
\end{proposition}
\begin{proof}
Connected cooriented components in a knot exterior have relative classes
$0,+1,-1$. Their sum is $+1$, so there are $q+1$ positive components.
Choose one. The union of all others has zero class. Its signed edge
intersection cocycle is exact, say $\delta g$, so deletion replaces the
coherent edge cochain by the remaining cochain in the same primitive class.
The remaining absolute edge weights uniquely determine its admissible
normal coordinates. In one tetrahedron the kernel of the six edge-weight
map is generated by $(1,1,1,1,-1,-1,-1)$; two distinct admissible
nonnegative vectors cannot differ by a nonzero multiple of it.
Thus deletion leaves a coherent representative and the chosen component
has at least $D_*$ pieces. Applying this also with orientation reversed
to each nonzero component gives $(2q+1)D_*$ nonzero-class mass.
Every component uses a piece, giving the count bound. Finally,
nonnegative edge weights give $\chi(S\setminus C)\le D(S\setminus C)$.
\end{proof}

The report's nine-tetrahedron $S^1$-times-triangle staircase attains the
opposite-pair bound for every $q$: $D_*=3$, with $q+1$ positive and $q$
negative three-piece components, hence $\delta=6q$. These are uniform
geometric statements, independent of the finite performance corpus.

\subsection{Why edge-first optimization produces one component}

The primary problem is the existing unconstrained integral min-$F$
network. Any independently verified optimal dual $f$ identifies its whole
optimal face edge by edge, with $x_e=c_e+p_b-p_a$:
\[
 \begin{array}{c|c}
 w_e=0&\text{no restriction}\\
 |f_e|<w_e&x_e=0\\
 f_e=w_e&x_e\ge0\\
 f_e=-w_e&x_e\le0.
 \end{array}
\]
This is complementary slackness for $w_e|x_e|$, not an extra choice of
height chamber. Add upper and lower variables $U_t,L_t$ with constraints
$p_{v(t,i)}-U_t\le-h_{t,i}$ and
$L_t-p_{v(t,i)}\le h_{t,i}$. Minimizing $\sum_t|U_t-L_t|$ on that
face is the secondary span network. The two integral networks take
$O(T^2\log T)$ indexed operations in the report's elementary model, with
separate polynomial bit costs; they do not expand the height range.

\begin{proposition}[Connected lexicographic optimum and success threshold]
Every nonempty primitive coherent lexicographic minimum of $(F,D)$ in
a knot exterior is connected and orientable. If $F_*\le2D_*$, it is an
essential disc or an annulus with exactly one inessential boundary circle,
whose cap gives an essential disc. In particular, a coherent essential
disc with at most $D_*+1$ pieces guarantees success.
\end{proposition}
\begin{proof}
Select a positive component as above and coherently delete the zero-class
remainder. If its edge penalty is positive, deletion lowers the primary
objective; otherwise a nonempty remainder lowers the secondary span.
Both contradict optimality. Thus the optimum is connected. Its Euler
characteristic is at least $D_*-F_*/2\ge0$. Its primitive relative class
forces nonempty boundary, so classification of connected orientable
surfaces leaves a disc or annulus. Two essential parallel annulus boundaries
have total class zero or twice a primitive class, neither the prescribed
primitive class; two inessential boundaries have class zero. Exactly one
boundary is therefore inessential in the annulus case. Finally a disc
with $D\le D_*+1$ has $F=2D-2\le2D_*$.
\end{proof}

Negative Euler characteristic is a failed candidate, not a knottedness
certificate. The primitive/source hypothesis is part of the theorem and
cannot be inferred from an optimizer's arithmetic output.

\subsection{Native proof formats and avoiding the second flow}

\begin{sloppypar}
The arithmetic producer \code{cocycle\_lex.py} emits
\code{normal-cocycle-edge-span-v1}: global vertex identifiers, a primary
primal/dual proof, a secondary proof, piece count and Euler characteristic.
The independent checker \code{cocycle\_lex\_verify.py} rebuilds the whole
optimal face from the actual source and verifies both objectives without
importing either optimizer. The outer \code{diagram-cocycle-lex-v1} binds
the result to a canonical input diagram, its authenticated exterior and
the primitive cochain. Actual normal edge weights, coordinates and Euler
characteristic are replayed. Boundary shelling transports this source proof
through the existing independently checked shelling record.
\end{sloppypar}

The maintained pipeline already has an independently verified global
minimum-span dual. If the primary optimum attains that $D_*$, it is
automatically also secondary optimal. We retarget the old span proof to
the new coordinates and check it again; no second network is needed.
Otherwise the producer runs the full secondary network. Both cases have
independent replay and the same lexicographic contract.

The opt-in stage runs after the four-tree prelude and optional minimum-span
stage, before additional tree trials. It reuses prepared source geometry.
Only a successful disc/cappable-annulus candidate receives the outer source
proof. Callback cancellation, errors and resource limits propagate through
the producer and checker. The existing recognizer fallback remains available.
The public switch is \code{normal\_seed\_edge\_span=True}, or CLI
\code{--normal-seed-edge-span}; the default remains false.

\subsection{Validation and the reason this stage remains optional}

The first native release passes 1,356 tests in 244.385 seconds. Its focused
37-test run includes one optional-dependency skip in the system interpreter;
the full native run includes Regina. The audit exactly preserves all 85
legacy default answers and work records when the option is disabled.
With it enabled, positives remain 24 of 85. Arithmetic completes on 84
sources, one is capped, 81 reuse the old span dual and three require the
second network. Nineteen source-bound positive proofs and 75 fresh Regina
connectedness/orientability/Euler/disc controls agree. On two random
examples Euler characteristic improves from $-3$ to $-1$, still no disc.
Mutation, 20,000-bit gauge, producer-disabled replay, shelling and
cancellation tests exercise both secondary formats.

Two serial four-arm benchmarks each complete 220 measured calls and 44
warmups, with full seed/recognition proofs, replay and serialization.
Unrelated Lean workers shared the host; both raw runs and A/A controls
are retained, and these measurements are not presented as isolated-host
timings. Early positive controls stay near parity. Failed seed queries
become slower, with old/new ratios about $0.42$--$0.74$; full recognition
miss controls have ratios about $0.70$--$0.74$. This is roughly
35--43 percent additional full-recognition time, not a speedup.

\begin{samepage}
\noindent\textbf{Full recognition controls from the second shared-host run.}
\begin{center}
\small
\begin{tabular}{lrrr}
\toprule
Complete recognition & default ms & enabled ms & default/enabled\\
\midrule
optimized positive & 50.490 & 50.988 & 0.993\\
genus one miss & 20.867 & 29.597 & 0.721\\
trefoil & 18.081 & 25.402 & 0.711\\
figure eight & 32.220 & 43.575 & 0.738\\
\bottomrule
\end{tabular}
\end{center}

\end{samepage}

No observed coverage gain justifies default activation. The new theorem
and independently replayable candidate remain useful for future adaptive
policies or sources where coherent span excess is demonstrably small.
The frozen baseline is \code{5402cd990}; evidence and the maintained
\code{normal\_orbit\_research.lex} driver are retained under \code{data/lex-*}.
